\subsection{Learning objectives and exact endpoint targets}\label{app:learning-objectives}
This identity suggests learning the reference's weighted demand response
where it affects spatial accounts. With training-only positive scales
$a_j$, the boundary term is
\begin{equation}
\mathcal L=\mathcal L_{\mathrm{flow}}+\lambda_J\mathcal L_{\mathrm{response}}
+\lambda_{\mathrm{cut}}\mathcal L_{\mathrm{boundary}},\qquad
\mathcal L_{\mathrm{boundary}}
=\frac1{NK}\sum_{q,j}\left[\frac{\mathcal K_j(h_j(q))}{a_j}\right]^2.
\label{eq:trainingobjective}
\end{equation}
The reference Jacobian is obtained from the existing implicit SUE
sensitivity system (Appendix~\ref{supp:sue}). We compare ordinary
flow supervision, flow plus response supervision, and the full objective.
Derivative supervision itself is established in Sobolev training
\citep{czarnecki2017sobolev}, including finite-difference derivative targets
\citep{yuan2022sobolev}; the network-derived target here selects the
part of weighted response error that changes boundary accounts.
A circulation in $\ker B$ is invisible to $\mathcal K_j$, so flow and
response losses remain necessary for the other attribution objectives.
The target matches the reference's actual response; it does not force
its true internal attribution divergence to zero or replace TSTT weights
with uniform weights. Equation~\eqref{eq:localboundarybound} is a path
inequality; finite training-point losses alone do not certify an entire
demand domain.


For the following finite-change statement, the outcome is explicitly separable:
\[
g(F)=(g_e(F_e))_{e\in E},\qquad T(F)=\sum_{e\in E}g_e(F_e).
\]
\begin{proposition}[exact finite-change boundary target]\label{prop:pairboundary}
Let $f^L,f^S$ and the edgewise outcome $g$ be continuously differentiable
on an open demand domain, and let $q^\pm=q\pm h e_j$ with their entire
segment in that domain. Put
\[
z_j^r=g(f^r(q^+))-g(f^r(q^-)),\qquad r\in\{L,S\}.
\]
For the attribution of this single-OD change,
\begin{equation}
\max_{U:s_j\in U,\ t_j\notin U}|\phi_j^L(U)-\phi_j^S(U)|
=\mathcal K_j(z_j^L-z_j^S).\label{eq:pairboundaryexact}
\end{equation}
If all sufficiently small single-coordinate changes have matching
accounts on every separating boundary at every point of the domain,
then $BD^LJ_j^L=BD^SJ_j^S$ throughout the domain. Integrating yields
matching boundary attribution columns for any joint demand path in that
domain (with the same columnwise path convention).
\end{proposition}
Appendix~\ref{app:pairproof} proves both statements. The first is an
exact finite-change target, whereas the pointwise response objective
optimizes the upper bound in Eq.~\eqref{eq:localboundarybound}; integrating
before taking $\mathcal K_j$ can retain cancellations. The second statement
requires a domain-wide local premise and is not a guarantee from finite
training data. With several ODs changing, the endpoint outcome difference
only identifies the sum of their attribution columns.

For sampled single-OD endpoint pairs, a direct training objective is
\begin{equation}
\mathcal L_{\mathrm{pair}}=\mathcal L_{\mathrm{flow}}
+\lambda_{\mathrm{edge}}\frac1N\sum_{(q,h,j)}
\frac{\|z_j^L-z_j^S\|_2^2}{b_j^2}
+\lambda_{\mathrm{cut}}\frac1N\sum_{(q,h,j)}
\left[\frac{\mathcal K_j(z_j^L-z_j^S)}{a_j}\right]^2,
\label{eq:pairobjective}
\end{equation}
with positive training-only scales $a_j,b_j$. The edge term controls
changes invisible to the boundary seminorm, whose kernel is $\ker B$.
Two reference endpoint equilibria supply $z_j^S$; training uses ordinary
parameter backpropagation through two predictions, without reference
Jacobian labels or input-Jacobian differentiation. Actual label-generation
and training costs must still be measured.
